\begin{definition}[Literal transports on the three selected vertical bonds]%
    \label{def:peps_torus_three_plaquette_flux_assignment}
    \lean{TNLean.PEPS.torusThreePlaquetteFluxAssignment,
        TNLean.PEPS.torusThreePlaquetteFluxAssignment_eq_treeCycleAssignment,
        TNLean.PEPS.torusThreePlaquetteFluxAssignment_up_transport}
    \leanok
    \uses{thm:peps_three_plaquette_geometry}
    In the actual translated three-plaquette strip, prescribe upward
    transports $t_0,t_1,t_2$ on its first three vertical bonds. The rightmost
    vertical bond and all other bonds carry the identity.
    In the native ordered convention, the three selected coefficients are
    $k_y(t_i)$, where
    \begin{align}
        k_y(t)=
        \begin{cases}
            t,      & \operatorname{val}(y)<\operatorname{val}(y+1), \\
            t^{-1}, & \text{otherwise}.
        \end{cases}
    \end{align}
    This literal assignment is exactly the tree-cycle assignment with
    residual $k_y(t_i)$ on $e_i$, and identity on every other non-tree bond.
    Thus a single orientation flag treats all three chords, including
    at the vertical seam. This construction is auxiliary to the adjacent
    fluxes in \cite[braiding passage]{Schuch2010PEPS}.
\end{definition}

\begin{theorem}[Actual three-plaquette and surrounding holonomies]%
    \label{thm:peps_torus_three_plaquette_flux_holonomy}
    \lean{TNLean.PEPS.regularWalkHolonomy_torusThreePlaquetteFluxAssignment,
        TNLean.PEPS.torusThreePlaquette_rightPairWalk_mem_region}
    \leanok
    \uses{def:peps_torus_three_plaquette_flux_assignment,
        def:peps_torus_plaquette_flux_geometry,
        def:peps_torus_joint_flux_geometry}
    The actual counterclockwise left, middle, and right plaquette walks
    have holonomies
    \begin{align}
        (a,b,c)=(t_0^{-1}t_1,\ t_1^{-1}t_2,\ t_2^{-1}).
    \end{align}
    The six-step outer walk surrounding the middle and right plaquettes
    has holonomy $t_1^{-1}=bc$. Every site visited by this walk belongs
    to the actual eight-site strip. These formulas hold at every translated
    position under horizontal period at least five and vertical period at
    least three, including both seams.
    They concern an explicit finite coordinate geometry; an identification
    of a corresponding physical operation with a string-crossing braid is
    not asserted.
\end{theorem}

\begin{proof}\leanok
    Directed upward transport at the four columns is respectively
    $t_0,t_1,t_2,1$. All horizontal transports are the identity, and reversal
    inverts the directed transport. Substitution in the reverse-ordered
    walk products gives the three plaquette formulas and the surrounding
    value $t_1^{-1}$. The six vertices of the latter walk occupy labels
    $1$ through $6$ of the eight-site numbering, proving containment.
\end{proof}
